\section{组织与人：反人性、顺人性与共同大脑}

上一章讲平台，本章回到组织。李想反复提到“顺着人性去说，逆着人性去做”。这句话容易被理解成管理技巧，但在访谈语境里，它指向组织变革的困难：最佳实践、能力建设、长期主义、复盘和自我否定往往反人性；而随心所欲、短期逃避、自我合理化更顺人性。一个组织如果要持续进化，就必须设计一种既理解人性、又能推动反人性成长的机制。

\begin{knowledgebox}{反人性不是压抑人，而是对抗惰性}
当李想说“反人性”，他强调的是对抗懒惰、逃避和短期舒适，而不是否定人的感受。优秀组织必须顺着人的表达和尊严去沟通，但在能力建设和结果要求上不能放任。
\end{knowledgebox}

\subsection{三到七个人的共同大脑和共同心脏}

上一节说明组织要理解人性但不能纵容惰性，本节看李想给出的具体组织单元。访谈后半程，他谈到一种模式：三到七个人组成更强的大脑、更强的心脏和更强的能量。这不是普通小组讨论，而是把能力互补、情绪能量、责任共担和方法论统一放在一起。一个人很有能力但没有结果，可能是因为缺少这种共同结构。

\begin{importantbox}{共同大脑的四个条件}
第一，成员彼此互补，而不是互相复制。第二，有共同目标和方法论。第三，关系足够近，能够真实反馈。第四，团队有共同承担压力的能量，而不只是分工表。
\end{importantbox}

\subsection{在意用户与在意身边的人}

李想把人和人连接的本质归纳为两个“在意”：在意用户，在意身边的人。在意用户是价值观共识，在意身边的人则是协作基础。一个组织如果只讲事情、不讲人，很容易把关系消耗掉；如果只讲关系、不讲用户，又会失去结果。优秀组织要在两者之间保持张力。

\begin{warningbox}{“对事不对人”有时会掩盖真实问题}
工程组织常说对事不对人，这能避免情绪化攻击。但如果完全不看人，就会忽略人的成长、动机、能量和关系质量。李想这里强调“先对人再做事”，是提醒组织问题最终要落到人是否变强。
\end{warningbox}

\subsection{本章小结}

组织能力不是流程图，而是人如何在长期压力中共同成长。AI 时代的小组织可能更强，但前提是它有共同大脑、共同心脏、共同方法论和共同能量。

